%!TEX root = hems_proj.tex
%%%%%%%%%%%%%%%%%%%%%%%%%%%%%%%%%%%%%%%%%%%%%%%%%%%%%%%%%%%%%%%%%%%%%%%
\chapter{Következtetések és További Célok}\label{ch:ZARSZO}
%%%%%%%%%%%%%%%%%%%%%%%%%%%%%%%%%%%%%%%%%%%%%%%%%%%%%%%%%%%%%%%%%%%%%%%

\begin{osszefoglal}
A fejezet összegzi a dolgozat fő megállapításait, kitér a gazdasági és környezeti vonatkozásokra, végül felvázolja a továbbfejlesztés lehetséges irányait.
\end{osszefoglal}

\section{Összegző következtetések}
A dolgozat fő tanulsága, hogy a HEMS-optimalizálásban a bemeneti adatok realizmusa legalább annyira meghatározó, mint maga az algoritmusválasztás. A szezonális és a 2x2-es kísérletek együtt azt mutatták, hogy a tisztán szintetikus környezetben elért teljesítmény nem vetíthető ki közvetlenül a valós üzemre: amint a termelési és a fogyasztási oldalon is valós, zajos idősorok jelennek meg, a keresési tér rücskössé válik, és az algoritmusok közötti különbségek felerősödnek. Ebben a megváltozott környezetben a populáció sokszínűségét fenntartó genetikus algoritmus, különösen a 100-as populációval futtatott változat bizonyult a legstabilabbnak, míg a gyors, de a globális legjobbhoz erősen vonzódó PSO a valós bemenetekre érzékenynek mutatkozott.

Fontos kiemelni, hogy nem létezik egyetlen, minden helyzetben legjobb módszer. A szűkös téli termelés mellett a kevés paramétert igénylő szürkefarkas-optimalizáló adta a legjobb átlagot, a bőséges nyári termelésnél a genetikus algoritmus finomhangolási előnye érvényesült, a hibrid eljárás pedig kiegyensúlyozott, de nem mindig kiemelkedő teljesítményt nyújtott. A választás ezért az üzemeltetési prioritástól függ: a stabilitás, a futási idő vagy a paraméterigény eltérő súlyozása más-más módszert tesz előnyössé. A projekt értékét nem utolsósorban az adja, hogy valós árinformációra épül, szétválasztja a bemeneti realizmus és az algoritmikus teljesítmény hatását, és a nyers adatok, valamint az összegző mutatók mentése révén minden megállapítása visszakövethető.

\section{Gazdasági és környezeti megfontolások}
A modell gyakorlati relevanciájának érzékeltetésére érdemes egy indikatív gazdasági becslést is megadni, hangsúlyozva, hogy ez nem mért eredmény, hanem a komponensek közelítő romániai piaci áraira és a szimulált napi költségcsökkenések éves szintre vetítésére épülő nagyságrendi tájékozódás. A \ref{tab:gazdasagi}. táblázat egy közepes méretű háztartás tipikus konfigurációját veszi alapul. Az itt szereplő éves megtakarítás a kísérletekben tapasztalt napi költségcsökkenésekből származó, évszakok közötti átlagolással becsült érték, ezért a tényleges megtérülés a fogyasztási szokásoktól, a tarifától és az időjárástól függően érdemben eltérhet.

\begin{table}[H]
\centering
\caption{Indikatív beruházási és megtérülési becslés egy közepes háztartásra}
\label{tab:gazdasagi}
\begin{tabular}{lc}
\toprule
Tétel & Közelítő érték \\
\midrule
HEMS-hardver & 1\,200 RON \\
5 kWp napelemes rendszer & 15\,000 RON \\
10 kWh akkumulátor & 18\,000 RON \\
\midrule
Teljes beruházás & 34\,200 RON \\
Becsült éves megtakarítás & 6\,200 RON \\
Becsült megtérülési idő & 5--6 év \\
\bottomrule
\end{tabular}
\end{table}

A megtérülés a romániai Casa Verde típusú támogatásokkal jellemzően három-négy évre rövidülhet. A környezeti hatás szintén számottevő, bár szintén becslés jellegű: egy évi mintegy 6\,000 kWh saját napelemes termelés nagyságrendileg 4--5 tonna szén-dioxid kiváltását jelentheti, amihez a terheléseltolásból és a hatékonyságnövelésből származó további megtakarítás társul. Ezek a számok a rendszer rendeltetésszerű, teljes kihasználását feltételezik, ezért inkább a lehetséges hatás nagyságrendjét, mintsem pontos értékét érzékeltetik.

\section{További fejlesztési irányok}
A jelenlegi modell több irányban is továbbfejleszthető, és ezek a lépések egyúttal az eredmények tudományos megalapozottságát is erősítenék. A legközvetlenebb feladat magának a kiértékelésnek a megerősítése: a jelenlegi, profilonként két futás helyett algoritmusonként és profilonként legalább tíz-húsz ismétlés, valamint az ezekre épülő konfidencia-intervallumok használata megbízhatóbbá tenné az összehasonlítást, és csökkentené az egy-egy szerencsés vagy kedvezőtlen seed torzító hatását. Ehhez kapcsolódóan érdemes volna egy egységes alapvonalat, akkumulátor és ütemezés nélküli esetet, valamint rögzített tanuló- és validációs naplistát is definiálni.

A modell oldaláról a legfontosabb pontosítás a célfüggvény bővítése: a jelenlegi költségmodell még nem tartalmaz explicit akkumulátor-degradációs költséget, a visszatáplálási árat pedig egyszerűsített szorzóval kezeli, így a gazdasági becslések finomításához ezeket a tényezőket érdemes lenne részletesebben modellezni. Tartalmi irányban a vizsgálat kiterjeszthető több, egymástól eltérő profilú háztartásra és speciális fogyasztókra, például hőszivattyúra vagy elektromos jármű töltésére, valamint hosszabb, akár éves, megszakítás nélküli időhorizontra.

Hosszabb távon több, a szakirodalomban is ígéretes irány érdemel figyelmet. Az energiaközösségek és a felhasználók közötti energiakereskedelem bevonása a megosztott tárolás révén tovább csökkentheti a tagok költségeit \cite{Lazzari2023}; a többcélú, Pareto-front alapú optimalizálás lehetővé tenné, hogy a felhasználó maga válasszon a költségminimalizálás és a komfort között \cite{Fadaee2012}; a gépi tanulás, például az időjárás- és fogyasztás-előrejelzés integrálása pontosabb, előretekintő ütemezést adhatna; a jármű-hálózat (Vehicle-to-Grid) megközelítés pedig az elektromos autók akkumulátorát mobil tárolóként vonná be a rendszerbe. Ezek az irányok együttesen mutatják, hogy a bemutatott keret nemcsak önmagában értékes, hanem szilárd alapot ad a további, valós környezetben is helytálló kutatáshoz.
